\section*{Capítulo \ref{ch:b2:cell-signalling} --- Mensageiros químicos e transdução de sinais}

\begin{solution}{exo:b2:cell-signalling:1}
Insulina: endócrina, hidrossolúvel (peptídeo). Cortisol: endócrino,
lipossolúvel. Acetilcolina numa sinapse: sináptica, hidrossolúvel.
Histamina numa ferida: parácrina, hidrossolúvel. Estradiol: endócrino,
lipossolúvel. Adrenalina da suprarrenal: endócrina, hidrossolúvel.
Óxido nítrico: parácrino, lipossolúvel (um gás que atravessa as
membranas).
\end{solution}

\begin{solution}{exo:b2:cell-signalling:2}
A adrenalina se liga ao receptor (desligamento: dissociação,
internalização do receptor); o receptor ativa a \omterm{prop:b2:cell-signalling:gpcr}{proteína G}, trocando GDP
por GTP (desligamento: hidrólise do GTP em segundos); a \omterm{prop:b2:cell-signalling:gpcr}{proteína G} ativa
a adenilil ciclase, que produz AMPc (desligamento: fosfodiesterase); o
AMPc ativa a quinase A (desligamento: perda do AMPc); a quinase A
fosforila a fosforilase quinase, que fosforila a glicogênio fosforilase
(desligamento: fosfatases); a fosforilase libera glicose-1-fosfato do
glicogênio, que é convertida e exportada como glicose.
\end{solution}

\begin{solution}{exo:b2:cell-signalling:3}
Os receptores de superfície ligam mensageiros hidrossolúveis na
membrana e agem em segundos por meio de \omterm{prop:b2:cell-signalling:gpcr}{segundos mensageiros} e de
quinases, mudando a atividade de proteínas já existentes. Os \omterm{prop:b2:cell-signalling:nuclear}{receptores
nucleares} ligam mensageiros lipossolúveis dentro da célula e agem em
horas como \omterm{prop:b2:cell-differentiation:transcriptional}{fatores de transcrição}, mudando quais proteínas são
produzidas.
\end{solution}

\begin{solution}{exo:b2:cell-signalling:4}
\omterm{prop:b2:cell-signalling:gpcr}{AMP cíclico}: produzido pela adenilil ciclase, destruído pela
fosfodiesterase. Cálcio: entra por canais ou sai do retículo
endoplasmático em resposta ao inositol trifosfato, removido por bombas.
Inositol trifosfato (com o diacilglicerol): produzido pela fosfolipase C
a partir de um lipídio de membrana, removido por fosfatases.
\end{solution}

\begin{solution}{exo:b2:cell-signalling:5}
$\theta = [L]/(K_d + [L])$: $0.09$, $0.5$, $0.91$, $0.99$. Dobrar os
receptores deixa a fração inalterada e dobra o número de receptores
ocupados --- e, portanto, a resposta.
\end{solution}

\begin{solution}{exo:b2:cell-signalling:6}
$50\times 200\times 100\times 300 = 3\times 10^{8}$. Cem receptores
ocupados dão $3\times 10^{10}$ moléculas de produto por segundo --- um
limite superior, pois as etapas saturam.
\end{solution}

\begin{solution}{exo:b2:cell-signalling:7}
A \omterm{prop:b2:cell-signalling:gpcr}{proteína G} travada mantém a ciclase ativa, o AMPc permanece alto, a
quinase A mantém fosforilado e aberto o canal de cloreto das células
intestinais, o cloreto é secretado, e o sódio e a água o seguem para a
luz --- litros por dia. A trava é uma modificação covalente da \omterm{prop:b2:cell-signalling:gpcr}{proteína
G}, desfeita apenas quando as próprias células são substituídas, dias
depois.
\end{solution}

\begin{solution}{exo:b2:cell-signalling:8}
$0.9\times 10^{-6}\,\text{mol/L}\times 2\times 10^{-12}\,\text{L} =
1.8\times 10^{-18}$ mol, $1.1\times 10^{6}$ íons (na verdade mais, pois
os tampões ligam a maior parte do que entra). Um único canal a $10^{6}$
por segundo levaria um segundo; mil canais o fazem em um milissegundo,
que é a escala de tempo de que um sinal precisa.
\end{solution}

\begin{solution}{exo:b2:cell-signalling:9}
As duas vias convergem sobre as mesmas enzimas: a quinase A do glucagon
fosforila a glicogênio sintase (desligando-a) e a fosforilase quinase
(ligando-a); a cascata da insulina ativa a fosfatase que remove os
fosfatos. O estado das enzimas é fixado pelo equilíbrio entre a
atividade das quinases e a das fosfatases, de modo que o que conta é a
razão entre os dois \omterm{def:b2:cell-signalling:kinds}{hormônios}; com ambos altos, a fosfatase em geral
vence e o glicogênio é armazenado.
\end{solution}

\begin{solution}{exo:b2:cell-signalling:10}
Adrenalina: o AMPc se difunde através de uma célula de
\qty{20}{\micro m} em menos de um segundo ($x^{2}/2D$ com
$D \approx \qty{3e-10}{m^2/s}$), e todas as enzimas da cascata já
existem --- segundos. Cortisol: um transcrito leva \qty{60}{s}, uma
proteína \qty{100}{s}; a partir de algumas dezenas de mRNAs, cada um com
uma dúzia de ribossomos, a célula produz algumas centenas de moléculas
de enzima por minuto, de modo que os milhares necessários para um efeito
mensurável levam horas, após os minutos de que o cortisol precisa para
entrar e chegar ao DNA.
\end{solution}

\begin{solution}{exo:b2:cell-signalling:11}
O receptor sinaliza continuamente: a cascata das MAP quinases leva a
célula a se dividir sem nenhum fator de crescimento, ignorando a
ausência do sinal que normalmente a contém. Como uma única \omterm{def:b2:genome-diversification:mutation}{mutação}
remove um controle sobre a divisão, esses receptores (e as quinases a
jusante) estão entre os oncogenes mais comuns. Um fármaco que ocupa o
sítio de ATP da quinase interrompe as fosforilações e silencia o
receptor --- o princípio de vários fármacos contra o câncer.
\end{solution}

\begin{solution}{exo:b2:cell-signalling:12}
Endócrino: chega a todas as células em um minuto, barato por célula,
não consegue se dirigir a uma única célula, age de minutos a dias ---
indispensável para os estados do corpo inteiro (jejum, crescimento,
reprodução). Nervoso: um milissegundo, um único alvo, fiação cara, age
durante milissegundos --- indispensável para o movimento e os reflexos
rápidos. Os dois são usados para a mesma mensagem quando se deseja ao
mesmo tempo velocidade e alcance: os nervos simpáticos e a medula
suprarrenal fornecem ambos noradrenalina ou adrenalina num susto.
\end{solution}

\begin{solution}{pb:b2:cell-signalling:1}
\textbf{1.} $1\,\text{nmol}/5\,\text{L} = \qty{0.2}{nmol/L}$.
\textbf{2.} $\theta = 0.2/1.2 = 0.17$: cerca de 3300 receptores
ocupados.
\textbf{3.} Três meias-vidas: \qty{0.025}{nmol/L}; $\theta = 0.024$.
\textbf{4.} Entre meio minuto e um minuto, o tempo que o \omterm{def:b2:blood-circulation:blood}{sangue} leva
para dar a volta; o nervo entrega seu transmissor diretamente ao fígado
em um segundo.
\textbf{5.} A mesma fração, mas $33\,000$ receptores ocupados: uma
resposta dez vezes maior.
\textbf{6.} A \qty{1}{nmol/L}, um receptor com $K_d = \qty{1}{\micro
mol/L}$ fica ocupado a um milésimo: nenhuma resposta. Ajustar o $K_d$ à
faixa circulante coloca a parte íngreme da curva de ocupação onde o
\omterm{def:b2:cell-signalling:kinds}{hormônio} de fato varia.
\textbf{7.} $20\times 100\times 1\times 50\times 50\times 1000 =
5\times 10^{9}$ moléculas de glicose por segundo por receptor ocupado.
\textbf{8.} $3300\times 20\times 100 = 6.6\times 10^{6}$ AMPc no
primeiro segundo; em \qty{5e-12}{L}, $1.1\times 10^{-17}$ mol, cerca de
\qty{2}{\micro mol/L}.
\textbf{9.} $3300\times 5\times 10^{9} = 1.7\times 10^{13}$ por segundo, $2\times
10^{13}$ por minuto por célula.
\textbf{10.} $2\times 10^{11}$ células $\times 10^{15} = 2\times
10^{26}$ moléculas por minuto, 330 mol, \qty{6e4}{g} por minuto ---
dez mil vezes os \qty{5}{g/min} observados. Os ganhos só se multiplicam
enquanto nenhuma etapa está saturada; na realidade, a fosforilase é
limitada por seu substrato e pelas fosfatases, e a exportação da glicose
pelos transportadores. O ganho da cascata é gasto em velocidade e
sensibilidade, e não em produção.
\textbf{11.} A \qty{5}{g/min}, vinte minutos; o susto termina em
minutos, e os mecanismos de desligamento param a cascata muito antes.
\textbf{12.} A hidrólise do GTP pela \omterm{prop:b2:cell-signalling:gpcr}{proteína G} (segundos); a
fosfodiesterase sobre o AMPc (segundos); as fosfatases sobre as enzimas
fosforiladas (de segundos a minutos); a \omterm{prop:b2:reproductive-hormones:pulses}{dessensibilização dos
receptores} (minutos). Os mais rápidos são o relógio próprio da \omterm{prop:b2:cell-signalling:gpcr}{proteína
G} e a fosfodiesterase.
\textbf{13.} O AMPc não é destruído: ele se acumula e persiste, a
quinase continua ativa, e a produção de glicose é maior e dura mais ---
o estado de alerta do café.
\textbf{14.} $4000/50 = \qty{80}{s}$.
\textbf{15.} $600/5 = \qty{120}{s}$ por enzima por ribossomo; 30 por
ribossomo por hora, 600 por mRNA por hora.
\textbf{16.} $100\times 600 = 60\,000$ enzimas por hora; $10^{5}$ após
cerca de \qty{1.7}{h}.
\textbf{17.} De dez a vinte minutos para o cortisol entrar, se ligar e
chegar a seus genes, uma hora até que os primeiros mRNAs sejam numerosos
e traduzidos, depois duas horas de acúmulo: cerca de quatro horas.
\textbf{18.} Uma cascata muda a atividade de enzimas que já existem; a
função do cortisol é mudar quais enzimas existem, durante dias, o que só
a transcrição pode fazer. Um susto não pode esperar quatro horas por
enzimas novas.
\textbf{19.} Mais receptores significam mais receptores ocupados em
qualquer concentração de adrenalina: a fração de ocupação não muda, mas
a resposta em cada concentração é maior --- a curva é ampliada, e não
deslocada. Um \omterm{def:b2:cell-signalling:kinds}{hormônio} lento fixa o ganho de um rápido.
\textbf{20.} A glicose entra na célula $\beta$ e é metabolizada; o ATP
sobe e fecha um canal de potássio sensível ao ATP; a membrana se
despolariza; canais de cálcio dependentes de voltagem se abrem; o
cálcio desencadeia a exocitose dos grânulos de insulina --- o metabolismo
acoplado à secreção por uma etapa elétrica.
\textbf{21.} Os dois receptores iniciam duas vias que convergem sobre as
mesmas enzimas: a quinase A do glucagon as fosforila, a cascata da
insulina ativa a fosfatase que as desfosforila; o estado das enzimas, e
portanto o sentido do metabolismo do glicogênio, segue a razão entre os
dois.
\textbf{22.} Glicemia de jejum alta (o fígado, sem freio, continua a
produzir glicose), glicemia pós-refeição muito alta (sem captação pelos
músculos e pelo tecido adiposo); sem insulina, as reservas de gordura
liberam ácidos graxos, que o fígado converte em cetoácidos ---
cetoacidose.
\textbf{23.} O músculo em atividade capta glicose sem insulina, de modo
que, numa corrida longa, a glicose cai enquanto a insulina injetada, sem
regulação, continua suprimindo o fígado: hipoglicemia e colapso.
\textbf{24.} Alça da glicose: sensores, as células $\beta$ e $\alpha$;
dois \omterm{def:b2:cell-signalling:kinds}{hormônios} efetores para os dois sentidos; um ganho alto, que mantém
o nível com precisão de um quinto; minutos para agir. \omterm{def:b2:blood-pressure:baroreflex}{Barorreflexo}:
receptores de estiramento, o \omterm{def:b2:heart:anatomy}{coração} e os vasos como efetores, ganho de
cerca de quatro, um segundo para agir.
\textbf{25.} Ocupação de $0.17$ após a liberação; ganho de $5\times 10^{9}$ por
segundo por receptor ocupado; produção do fígado de \qty{5}{g/min} na
realidade; o cortisol leva cerca de quatro horas.
\end{solution}
